% !TeX root = ../papers/nc-p3.tex
\section{二次剰余の全縦切片と二変数単元表示の有理化}
\label{sec:rank-fifteen-quadratic-verticals}
\begin{theorem}\label{thm:rank-fifteen-all-quadratic-verticals}
剰余標数3の付値で $\bar x\in\mathbb F_9$ なら、全ての非零の幾何水準 $y$ について、
通常な標数零の階数15窓は標準Koszul左右AS正則で、各始点が四つの中心四次式上自由階数256である。
特に3の上に剰余次数2以下の素点を持つ数体に $x$ が属すれば、任意の非零の $y$ にこの結論が成り立つ。
\end{theorem}
\begin{proof}
負付値の $x$ と素体に落ちる $x$ は既知である。残る六つの $X=\bar x$ は
$i^2=-1$ として $X=s+i$ ($s=0,1,-1$) とそのFrobenius共役である。
正付値の $y$ は単元横座標の内側の定理、負付値は極の定理で覆われる。
以下 $Y=\bar y$ は非零である。$c^2=X$, $b=1/c$, $d=b^3-b$, $D=v^3-v$ とし、
$Z^2=D^2-d^2$ と置く。前節の六根の積には
$W=(-ic)^3ZD^2/v^6$ が使える。実際、前節の六平方の積がこの二乗と一致する。
$Y=X^3D^2$ とし、素体上の写像で例外を扱うため $B=-D^2$ を用いる。
$B\in\mathbb F_9$ なら $Y\in\mathbb F_9$ なので単元同士の全二次剰余の定理を適用できる。

\path{computations/verify_rank_fifteen_quadratic_verticals.py} は36成分の表示で、
三つの $X$ とも前節と同じ未切断経路、文字0に五次元の四次中心、同じ自由列の後四本の切断を計算する。
四始点の次元は十三次の消滅まで標準商に一致する。
下表の四数は、入力と中心を含めた始点別の除外多項式の次数である。
\[
\begin{array}{c|c|r|r}
 X & \text{四始点} & \deg\mathrm{lcm} & \text{六表示で安定な部分の次数}\\ \hline
 i &(1205,1997,1898,2147)&6604&126\\
 1+i &(3240,3964,3281,2317)&12491&54\\
 -1+i &(3321,3604,3290,2843)&12756&48
\end{array}
\]
第一行では $c=1-i$, $d^2=-1$ なので前節の種数2の関数体を再利用し、
$B=0,1$ を除く。第二、三行では $c\in\mathbb F_{81}$ を使い、
各逆元のノルムを $\mathbb F_{81}(v)$ の多項式として記録する。
$B\in\mathbb F_9$ を除いた後、六変換 $v\mapsto\pm v+j$ の引き戻しの最大公約元を取る。
これと係数の三乗による四つの共役全体の最小公倍式は、素体上で次数72,48である。
$B$ の最小多項式ごとに完全ファイバーが得られることは前節と同じ引き戻しで照合する。

第一行で残る閉点の次数は $2,2,3,4,4,6$、第二、三行で素体上に拡大した除外集合は
それぞれ四次の三軌道と二軌道である。各素体多項式の根を九乗軌道に分け、
$Y=-X^3B$ ごとに元の七つの倍率を復元する。必要な体の素体上の次数は
第一行で $8,12$、第二行で $8,16$、第三行で8で足りる。各九乗軌道で四始点の十三次消滅までを
検査する。三乗は $X$ の他方の共役の全水準も与える。
安定除外集合の外ではいずれかの $v$ の全主元が単元になる。
四中心の持上げと自由加群比較は前節による。最後に $x$ が与えられた数体で負付値なら
第\ref{sec:rank-fifteen-negative-x-fan}節を、整なら上の結果を係数体への延長付値で使う。
\end{proof}

\begin{proposition}\label{prop:rank-fifteen-rational-unit-surface}
$X$ も動かすとき、$X\ne0,1$ かつ $Y\ne0,(1-X)^2$ の種数2表示は、
二変数の有理関数体による純非分離次数3の表示を持つ。
\end{proposition}
\begin{proof}
$S=v-b$, $T=v+b$ とすると、$F(z)=z^3-z$ の加法性から
$Z^2=F(S)F(T)$ である。この曲面は楕円曲線 $u^2=F(S)$ とそのコピーの積を
同時符号反転で割ったものに双有理である。$w=Z/F(S)$ の三乗根を $u$ とし、
$q=T-u^2S$ と取る。$F(T)=u^6F(S)$ を解けば
\[
 S=(q-q^3)/(u^6-u^2),\qquad T=q+u^2 S,\qquad Z=u^3 F(S)
\]
であり、$b=(T-S)/2$, $v=(S+T)/2$、従って四窓の全係数も戻る。
$F(S)F(T)=0$ は分岐水準である。$u\ne0$ で $u^4=1$ なら
$F(T)=\pm F(S)$ となり、$d=0$ または $D=0$ であって仮定に反する。
従ってこれらの式は指定した全幾何点で定義できる。
$S=0$ の上の一般分岐点で $F(S)$ の位数は2、$Z$ の位数は1である。
$w$ の位数は $-1$ で三乗ではないので、得た関数体の拡大次数はちょうど3である。
残る分岐水準では $F(S)=0$ または $F(T)=0$ となり、別の成分の規格化が必要である。
\end{proof}
